\bookchapter{Capítulo 27: Tus padres también pueden usarlo fácilmente - Guía completa de adaptación para mayores}

\chapterintro{Sinceramente, al escribir este capítulo pensé en mis padres. Mi madre, cada vez que el teléfono le muestra una “actualización del sistema”, no se atreve a pulsar, por miedo a romperlo. La fuente de WeChat de mi padre era más pequeña que una hormiga, pero no sabía cómo agrandarla y estuvo entrecerrando los ojos durante dos años. No es que sean “torpes”, es que los teléfonos por defecto tratan a todos como si tuvieran veintitantos años. En este capítulo, con palabras muy directas, resuelvo uno por uno los problemas más comunes que enfrentan los mayores al usar el teléfono.}

\booksection{Primero, ajusta tu mentalidad: no es que “los mayores no sepan usar”, es que “el teléfono no está bien diseñado”}

Antes de hacer cualquier ajuste, hay que dejar claro un concepto:\textbf{Las interfaces de los teléfonos inteligentes actuales no han sido diseñadas teniendo en cuenta a los usuarios mayores}. Las ventanas emergentes son tan pequeñas que no se ven, los gestos requieren toques precisos, los iconos se parecen todos: estos no son problemas del usuario, son problemas del producto. Así que ayudar a los mayores con el teléfono no consiste primero en enseñar, sino en ayudarles a\textbf{ajustar el entorno}. Si el entorno está bien ajustado, la mitad de los problemas desaparecen.

Hay tres reglas de oro para ayudar a los mayores con el teléfono: 1)\textbf{No les apresures}; tú llevas diez años usando smartphones, ellos pueden llevar solo tres meses. 2)\textbf{Usa analogías}; “el Wi-Fi es como las tuberías de agua de casa, abre el grifo y hay agua”, este tipo de metáfora es infinitamente más comprensible que la jerga técnica. 3)\textbf{Escríbelo}; si se lo enseñas de palabra pueden olvidarlo, pero si coges un papel y escribes los pasos clave y se lo pegas en la parte trasera de la funda del teléfono, funciona mejor que cualquier otra cosa.

\booksection{I. Primer paso: activa el “modo simple” - el teléfono se vuelve más claro}

Los teléfonos actuales suelen tener integrado un “modo simple” (llamado “modo cuidado”, “modo senior” o “modo mayor”), al activarlo, todo el teléfono cambia radicalmente:

\begin{itemize}
  \item \textbf{La fuente y los iconos se vuelven muy grandes}; adiós a las gafas para mirar el teléfono.
  \item \textbf{Los iconos del escritorio se simplifican drásticamente}; solo se mantienen teléfono, SMS, cámara, álbum, WeChat y otras appscomún, todo lo demás se guarda.
  \item \textbf{Las operaciones se simplifican}; muchas funciones vistosas (pantalla de la izquierda, asistente inteligente, ventanas flotantes, etc.) se desactivan por defecto, la interfaz se centra en la claridad.
  \item \textbf{El sonido aumenta}; el tono de llamada, el volumen de notificaciones y el volumen de llamada se ajustan automáticamente al máximo.
\end{itemize}

\booksubsection{Métodos de activación del modo simple por marca (sigue los pasos directamente):}

\begin{center}
\begin{tabularx}{\linewidth}{|X|X|X|}
\hline
\textbf{Marca} & \textbf{Nombre del modo} & \textbf{Método de activación} \tabularnewline
\hline
Sinceramente, es normal que el móvil empiece a irse lento tras un ano de uso. La causa no necesariamente es que este roto; más a menudo es "acumulacion de basura": archivos basura, procesos en segúndo plano, espacio de almacenamiento llenandose poco a poco. & Accesibilidad (sin modo simple separado, pero se puede ajustar individualmente) & Ajustes > Pantalla y brillo > Tamaño del texto, arrastra al máximo. Luego entra en "Accesibilidad > Pantalla y tamaño del texto" y activa "Texto en negrita" y "Forma de botones" \tabularnewline
Tengo un amigo al que le robaron el móvil el año pasado. Su primera preocupación no fue "cuánto vale el móvil", sino "¿podré recuperar mis fotos?". Lamentablemente, no había activado la copia de seguridad en la nube. El móvil es un objeto material, pero los datos de tu cuenta son cosas que, si se pierden, no vuelven nunca. En este capítulo hablaremos de cómo asegurar tu vida digital de verdad. & Modo simple & Ajustes > Sistema y actualizaciones > Modo simple > Activar \tabularnewline
Xiaomi/Redmi & Modo minimalista & Ajustes > Más ajustes > Modo minimalista > Activar \tabularnewline
OPPO/OnePlus & Modo simple & Ajustes > Funciones especiales > Modo simple \tabularnewline
vivo/iQOO & Escritorio simple & Ajustes > Escritorio, pantalla de bloqueo y fondo de pantalla > Escritorio simple \tabularnewline
Samsung & Modo simple & Ajustes > Pantalla > Modo simple \tabularnewline
\hline
\end{tabularx}
\end{center}

Si usas otra marca, busca “simple” o “mayor” en el cuadro de búsqueda en la parte superior de los “Ajustes” del teléfono; normalmente lo encontrarás. Si no lo encuentras, sigue los ajustes de “agrandar manualmente” a continuación uno por uno; el efecto es igual de bueno.

\booksection{II. Si no activas el modo simple, también puedes ajustar manualmente - “fuente grande + iconos grandes” en 5 pasos}

Algunos mayores no quieren usar el modo simple porque les parece demasiado “para mayores”. Puedes ajustar manualmente los siguientes cinco elementos con el mismo buen resultado:

\begin{enumerate}
  \item \textbf{Agrandar la fuente al máximo}: Ajustes > Pantalla > Tamaño de fuente, arrastra al extremo derecho. Aprovecha para ajustar también el "tamaño de visualización" (a veces llamado "escala de visualización") al máximo para que todos los elementos de la interfaz se agranden.
  \item \textbf{Iconos más grandes}: En la pantalla de inicio, junta dos dedos (como si reducieras una foto) para entrar en los ajustes del escritorio > Diseño del escritorio (o tamaño de iconos), elige el diseño con menos cuadrículas (como 3×5) y los iconos serán naturalmente más grandes.
  \item \textbf{Activar "texto de alto contraste"}: Hace los bordes del texto más nítidos; se puede encontrar en "Accesibilidad".
  \item \textbf{Desactivar notificaciones innecesarias}: Ajustes > Notificaciones; desactiva las notificaciones de todas las apps excepto teléfono, SMS y WeChat. Así el teléfono no sonará constantemente y no resultará molesto.
  \item \textbf{Poner contactos frecuentes en el escritorio}: Mantén pulsado un espacio vacío del escritorio > Añadir widget > Contactos; pon los accesos directos de llamada para tu pareja, hijos, médico y contactos de emergencia en el escritorio. Un toque y llama, sin buscar en los contactos.
\end{enumerate}

\begin{tcbpractice}{Práctica: ayuda a los mayores a "decorar" su teléfono}
\begin{enumerate}
  \item Activa el modo simple (sigue la tabla de marcas anterior).
  \item Activa la visualización del porcentaje de batería (Ajustes > Batería > Porcentaje de batería) para que siempre sepan cuánta batería queda.
  \item Despliega la barra de notificaciones y elimina u oculta los interruptores rápidos de Bluetooth, NFC, modo avión que los mayores no necesitan.
  \item Pon cuatro accesos directos de llamada en el escritorio: pareja, hijos, médico de familia (si lo hay), comunidad/comunidad de vecinos.
\end{enumerate}

\end{tcbpractice}

\booksection{III. Prevención de fraude - este capítulo debes leerlo con tus mayores}

Los ancianos son las principales víctimas del fraude telefónico. No es que sean poco cautelosos, sino que las tácticas de los estafadores están cuidadosamente diseñadas para explotar las debilidades psicológicas de los mayores. Estas son las estafas más comunes y cómo afrontarlas. Te recomiendo que se las leas o que imprimas esto y lo pegues junto al teléfono.

\booksubsection{Seis estafas comunes (si ves esto, cuelga inmediatamente):}

\begin{itemize}
  \item \textbf{1) "Tu nieto ha tenido un accidente, envía dinero rápido"}; el estafador dirá "tu nieto se ha metido en problemas, la policía lo ha detenido, necesita una fianza". Recuerda:\textbf{cualquier llamada que pida una transferencia para "soltar a alguien" es una estafa}. La policía no pide dinero por teléfono.
  \item \textbf{2) "Tu tarjeta sanitaria ha sido suspendida"}; un SMS dice que la tarjeta sanitaria tiene un problema y te pide hacer clic en un enlace para "activarla". Recuerda:\textbf{la seguridad social y el seguro médico no envían SMS con enlaces}. Si te ocurre esto, llama directamente al 12333 (línea oficial de seguridad social) para verificar.
  \item \textbf{3) "Inversión de alto rendimiento, sin riesgo"}; en grupos de WeChat o por SMS recomiendan "productos financieros para jubilados" o "proyectos de apoyo estatal" prometiendo altos rendimientos. Recuerda:\textbf{todo lo que prometa más del 5\% de rentabilidad anual y "garantice el capital" es una estafa en el 99,9\% de los casos}.
  \item \textbf{4) "Soy tu jefe/viejo amigo, he cambiado de número"}; se hace pasar por un conocido pidiendo dinero. Recuerda:\textbf{llamar para verificar es la mejor defensa}; si dicen "no puedo contestar el teléfono", el 80\% de las veces son estafadores.
  \item \textbf{5) "Huevos gratis/conferencia de productos sanitarios"}; actividades gratuitas presenciales que atraen a los ancianos a una "conferencia de salud" y luego les venden productos sanitarios a precios astronómicos. Recuerda:\textbf{no hay nada gratis en la vida, detrás de las actividades gratuitas siempre hay un negocio}.
  \item \textbf{6) "Tu teléfono tiene un virus, necesita reparación remota"}; el estafador te pide descargar una app de control remoto (como Sunlogin, TeamViewer) y luego controla tu teléfono para transferir dinero. Recuerda:\textbf{¡Nunca dejes que nadie controle tu teléfono remotamente!}Los fabricantes de teléfonos, operadores y bancos nunca te llamarán para pedirte que instales software de control remoto.
\end{itemize}

\booksubsection{Tres líneas de defensa; una vez establecidas, pueden bloquear el 90\% de las estafas:}

\begin{enumerate}
  \item \textbf{Primera: bloqueo integrado del teléfono}; en Android activa el "bloqueo de spam" (en la app de teléfono o en la app de gestión del teléfono); en iPhone activa "silenciar llamadas desconocidas" (Ajustes > Teléfono > Silenciar llamadas desconocidas).
  \item \textbf{Segunda: bloqueo del operador}; envía "KTFSR" al 10086 para activar gratis el bloqueo de llamadas de spam de China Mobile; China Telecom y China Unicom tienen servicios similares. Llama al operador y pide que lo activen.
  \item \textbf{Tercera: app del Centro Nacional Antifraude}; descarga la app "Centro Nacional Antifraude" de la tienda de aplicaciones, regístrate y activa la alerta de llamadas. Esta app está conectada con la base de datos de la policía y puede identificar los números de estafadores al instante.
\end{enumerate}

\begin{tcbwarn}{"Eslogan antifraude" para los mayores (imprímelo y pégalo junto al teléfono)}
\begin{itemize}
  \item \textbf{Llamada desconocida pidiendo transferencia ¡Cuelga!}
  \item \textbf{SMS con enlace para hacer clic ¡Bórralo!}
  \item \textbf{Prometen alto rendimiento garantizado ¡Bloquéalo!}
  \item \textbf{Te piden instalar software de control remoto ¡Cuelga directamente!}
  \item \textbf{¡La policía real no tramita casos por teléfono ni pide dinero!}
  \item \textbf{Si tienes dudas, llama primero a tus hijos para preguntar.}
\end{itemize}

\end{tcbwarn}

\booksection{IV. Asistencia remota - puedes ayudar a tus padres con el teléfono aunque estés lejos}

A veces los mayores tienen problemas y, por mucho que intentes describirlo por teléfono, no encuentran el botón. Aquí es donde la “asistencia remota” resulta útil: puedes ver la pantalla de su teléfono directamente desde el tuyo y ayudarles a operar. Cada marca tiene una solución diferente:

\begin{itemize}
  \item \textbf{Tengo un amigo al que le robaron el móvil el año pasado. Su primera preocupación no fue "cuánto vale el móvil", sino "¿podré recuperar mis fotos?". Lamentablemente, no había activado la copia de seguridad en la nube. El móvil es un objeto material, pero los datos de tu cuenta son cosas que, si se pierden, no vuelven nunca. En este capítulo hablaremos de cómo asegurar tu vida digital de verdad.}: La app integrada "Changlian" tiene función de "compartir pantalla". Tú y tus padres instaláis Changlian, iniciáis una videollamada y luego pulsáis "compartir pantalla"; podrás dibujar círculos en su pantalla para indicarles dónde tocar.
  \item \textbf{Xiaomi/Redmi}: La app de videollamada de Xiaomi tiene "asistencia remota" que soporta compartir pantalla + llamadas de voz simultáneas.
  \item \textbf{: puedes configurar Nginx o Apache para montar servidores web estáticos.}: Tiene "espacio familiar" o "guardián remoto" integrados; busca "remoto" en los ajustes para encontrarlos.
  \item \textbf{Solución universal (todas las marcas)}: Descarga la app "Sunlogin Remote Control" (la versión gratuita es suficiente); cada uno la instala en su teléfono. Registra la misma cuenta o añádelos como amigos; pide a tus padres que te dicten el código de verización que aparece en su teléfono y podrás ver y operar su pantalla remotamente.
  \item \textbf{Usuaríos de iPhone}: FaceTime soporta "compartir pantalla" durante las llamadas. En una llamada de FaceTime, toca el botón de compartir pantalla y la otra persona podrá ver tu pantalla. Pero las restricciones de seguridad de iOS no permiten operación remota; solo puedes "enseñar guiando".
\end{itemize}

\booksection{V. “Adaptación de WeChat para mayores” - con estas frases es suficiente}

WeChat es la app que más usan los mayores, pero la interfaz es demasiado densa. Estos son unos ajustes con efecto inmediato:

\begin{itemize}
  \item \textbf{Agrandar la fuente de WeChat}: WeChat > Yo > Ajustes > General > Tamaño de fuente, arrastra al máximo. Ten en cuenta que esto solo afecta a la fuente dentro de WeChat, no a la fuente del sistema.
  \item \textbf{Activar "modo cuidado"}: WeChat > Yo > Ajustes > Modo cuidado > Activar. Al activarlo, la fuente y los botones se agrandan, el contraste de color se refuerza y puedes "escuchar mensajes de texto" con un toque (los mensajes de texto se leen en voz alta).
  \item \textbf{Etiquetar contactos importantes para ellos}: Cambia los nombres de tus hijos, familiares en los contactos por algo fácil de reconocer, como "Hijo mayor - Xiao Ming" o "Hija - Xiao Hong", para que al abrir WeChat sepan quién es quién.
  \item \textbf{Enséñales a usar "entrada de voz" en lugar de teclear}: En el teclado junto al chat, mantén pulsada la barra espaciadora (Android) o toca el icono del micrófono (iPhone) y habla directamente; tu voz se convierte en texto y se envía. La precisión del reconocimiento ya es muy alta; para los mayores que no saben escribir en pinyin, es una comodidad "transformadora".
  \item \textbf{Desactivar funciones de "Descubrir" que no usan}: WeChat > Yo > Ajustes > General > Gestión de Descubrir; desactiva "Momentos", "Canales de video", "Directos", "Compras" y "Juegos" si no los usan, para reducir distracciones.
\end{itemize}

\booksection{VI. ¿Qué hacer si se pierde el teléfono? Prepara un “plan de emergencia” para tus mayores}

Mejor prevenir que curar. Antes de que se pierda el teléfono, haz estas cosas:

\begin{enumerate}
  \item \textbf{Confirma que "Buscar mi móvil" está activado}: iPhone en Ajustes > tu nombre > Buscar > Buscar mi iPhone; Android busca "buscar dispositivo" en los ajustes. Asegúrate de que este interruptor esté encendido.
  \item \textbf{Anota las contraseñas y cuentas importantes en un papel}: La contraseña de Apple ID / cuenta de Huawei / cuenta de Xiaomi; guárdala en un cajón cerrado y entrega una copia a un hijo de confianza. Si se pierde el teléfono, necesitarás acceder a estas cuentas para bloquear y localizar remotamente; si se te olvida la contraseña, será un problema.
  \item \textbf{Enséñales una frase a los mayores}:\textbf{"Si pierdo el teléfono, lo primero es llamar a la familia."}La familia puede iniciar sesión inmediatamente en iCloud o en la web de búsqueda de dispositivos de Android para marcarlo remotamente como "modo perdido"; así quien lo encuentre no podrá abrirlo ni usarlo.
\end{enumerate}

\booksection{VII. Auriculares inalámbricos y compatibilidad con audífonos - una buena función ignorada}

Muchos mayores tienen problemas de audición, pero no saben que el teléfono puede funcionar junto con audífonos. El iPhone y algunos Android soportan audífonos MFi (Made for iPhone), pudiendo enviar el sonido de la llamada directamente al audífono, mucho más claro que sostener el teléfono.

\begin{itemize}
  \item \textbf{Método de conexión}: iPhone Ajustes > Accesibilidad > Dispositivos de audición; activa Bluetooth y sigue las instrucciones de emparejamiento del audífono.
  \item \textbf{Si no usa audífono}: En "Accesibilidad > Audio/Visual", activa "Cancelación de ruido de llamada" y "Audio mono" para mejorar la claridad de las llamadas.
  \item \textbf{Android}: En "Accesibilidad > Mejora de audición" hay ajustes similares; la ubicación varía según la marca, busca "audífono" directamente.
\end{itemize}

\booksection{VIII. Funciones relacionadas con la salud - el teléfono no es solo una herramienta de comunicación}

\begin{itemize}
  \item \textbf{SOS de emergencia}: Pulsar el botón de encendido cinco veces seguidas (en la mayoría de teléfonos); el teléfono llama automáticamente al número de emergencia y envía la ubicación a los contactos de emergencia preestablecidos. Configúralo con antelación para tus mayores - iPhone: "Ajustes > SOS de emergencia"; Android: "Ajustes > Seguridad > Solicitud de emergencia". Enséñales a usarlo con un ensayo previo para que no se confundan en el momento.
  \item \textbf{Detección de caídas}: Apple Watch tiene esta función; si detecta una caída seguida de inactividad, llama automáticamente a emergencias y notifica a la familia. Si el presupuesto lo permite, compra uno para tus mayores; puede salvar vidas en momentos críticos.
  \item \textbf{Recordatorio de medicación}: La app "Salud" del iPhone tiene una función integrada de recordatorio de medicación; configura el nombre del medicamento, la hora y la dosis, y te recordará cuando llegue el momento. En Android puedes usar "alarmas programadas" o descargar apps de "recordatorio de medicación".
\end{itemize}

\begin{tcbsummary}{Resumen del capítulo}
\begin{itemize}
  \item Primero activa el modo simple o ajusta manualmente la fuente/iconos grandes; con el entorno bien ajustado, la mitad de los problemas desaparecen
  \item Tres reglas de oro para ayudar a los mayores con el teléfono: no apresures, usa analogías, escríbelo
  \item Seis estafas comunes que hay que recordar constantemente; tres líneas de defensa bien establecidas pueden bloquear el 90\% del fraude
  \item La asistencia remota te permite ayudar a tus padres con el teléfono aunque estés lejos
  \item El "modo cuidado" de WeChat + "entrada de voz" son las dos funciones más útiles para los mayores
  \item El SOS de emergencia y la detección de caídas son configuraciones que pueden salvar vidas y merecen 5 minutos de configuración
\end{itemize}

\end{tcbsummary}
